\begin{proposition}
\label{proposition-going-down-normal-integral}
Let $R \subset S$ be an inclusion of domains.
Assume $R$ is normal and $S$ integral over $R$.
Let $\mathfrak p \subset \mathfrak p' \subset R$
be primes. Let $\mathfrak q'$ be a prime of $S$
with $\mathfrak p' = R \cap \mathfrak q'$.
Then there exists a prime $\mathfrak q$
with $\mathfrak q \subset \mathfrak q'$
such that $\mathfrak p = R \cap \mathfrak q$. In other words:
the going down property holds for $R \to S$, see
Definition \ref{definition-going-up-down}.
\end{proposition}

\begin{proof}
Let $\mathfrak p$, $\mathfrak p'$ and $\mathfrak q'$
be as in the statement. We have to show there is a prime
$\mathfrak q$, with $\mathfrak q \subset \mathfrak q'$ and
$R \cap \mathfrak q = \mathfrak p$. This is the same
as finding a prime of
$S_{\mathfrak q'}$ mapping to $\mathfrak p$.
According to Lemma \ref{lemma-in-image} we have to show
that $\mathfrak p S_{\mathfrak q'} \cap R
= \mathfrak p$. Pick $z \in \mathfrak p S_{\mathfrak q'} \cap R$.
We may write $z = y/g$ with $y \in \mathfrak pS$ and
$g \in S$, $g \not\in \mathfrak q'$. Written
differently we have $zg = y$.

\medskip\noindent
By Lemma \ref{lemma-integral-integral-over-ideal}
there exists a monic polynomial
$P = x^m + b_{m-1} x^{m-1} + \ldots + b_0$
with $b_i \in \mathfrak p$ such that $P(y) = 0$.

\medskip\noindent
By Lemma \ref{lemma-minimal-polynomial-normal-domain}
the minimal polynomial of $g$ over $K$ has coefficients
in $R$. Write it as $Q = x^n + a_{n-1} x^{n-1} + \ldots
+ a_0$. Note that not all $a_i$, $i = n-1, \ldots, 0$
are in $\mathfrak p$ since that would imply
$g^n = \sum_{j < n} a_j g^j \in \mathfrak pS
\subset \mathfrak p'S
\subset \mathfrak q'$
which is a contradiction.

\medskip\noindent
Since $y = zg$ we see immediately from the above
that $Q' = x^n + za_{n-1} x^{n-1} + \ldots + z^{n}a_0$
is the minimal polynomial for $y$. Hence
$Q'$ divides $P$ and by Lemma \ref{lemma-polynomials-divide}
we see that $z^ja_{n - j} \in \sqrt{(b_0, \ldots, b_{m-1})}
\subset \mathfrak p$, $j =  1, \ldots, n$.
Because not all $a_i$, $i = n-1, \ldots, 0$
are in $\mathfrak p$ we conclude $z \in \mathfrak p$
as desired.
\end{proof}